\section{CTI Workflows}

Cyber Threat Intelligence becomes most valuable when it directly supports operational teams. Effective CTI workflows 
bridge the gap between intelligence reporting and hands-on defensive work by transforming threat insights into practical 
actions. For a CTI analyst, threat hunter, or detection engineer, the goal is not simply to collect intelligence but to 
operationalize it—turning adversary behaviors into hypotheses, detections, dashboards, and investigative leads. This 
chapter outlines the core workflows that enable CTI to drive meaningful security outcomes.

% --------------------------------------------------------------------
% SECTION 1: Identifying Relevant Intelligence
% --------------------------------------------------------------------

\subsection*{Identifying Relevant Intelligence}

A CTI workflow begins with identifying intelligence that matters to the organization. This may come from:

\begin{itemize}
    \item threat reports
    \item malware analysis
    \item internal telemetry
    \item external intelligence feeds
\end{itemize}

The analyst evaluates the intelligence to determine which behaviors, techniques, or infrastructure are relevant. 
Instead of focusing on short-lived indicators such as IP addresses or domains, modern CTI emphasizes techniques—how 
the adversary gains access, executes code, moves laterally, or evades detection. These behaviors form the foundation 
for operational work.

% --------------------------------------------------------------------
% SECTION 2: Translating Intelligence into Operational Questions
% --------------------------------------------------------------------

\subsection*{Translating Intelligence into Operational Questions}

Once relevant behaviors are identified, the next step is translation. CTI analysts convert high-level intelligence 
into concrete questions that can be answered with telemetry.

For example:

\begin{itemize}
    \item Intelligence: “Threat actor uses encoded PowerShell commands.”
    \item Operational question: “Where in our environment do we see PowerShell with encoded commands?”
\end{itemize}

This translation step is critical because it connects intelligence to the data sources available in Splunk or Sentinel.

% --------------------------------------------------------------------
% SECTION 3: CTI-Driven Threat Hunting
% --------------------------------------------------------------------

\subsection*{CTI-Driven Threat Hunting}

The translated intelligence becomes the basis for threat hunting. Analysts use SPL or KQL to search for behaviors 
associated with the threat actor.

\textbf{SPL Example}

\begin{lstlisting}
index=sysmon EventCode=1 Image="*powershell.exe"
| search CommandLine="*EncodedCommand*"
| table _time, Computer, User, ParentImage, Image, CommandLine
\end{lstlisting}

\textbf{KQL Equivalent}

\begin{lstlisting}
Sysmon
| where EventID == 1 and Image endswith "powershell.exe"
| where CommandLine contains "EncodedCommand"
| project TimeGenerated, Computer, User, ParentImage, Image, CommandLine
\end{lstlisting}

If the hunt reveals suspicious activity, the workflow moves into detection engineering.

% --------------------------------------------------------------------
% SECTION 4: CTI-Driven Detection Engineering
% --------------------------------------------------------------------

\subsection*{CTI-Driven Detection Engineering}

When CTI-driven hunts uncover suspicious behavior, the next step is to refine the behavioral pattern into a repeatable 
detection. This ensures that the organization is protected not only against the current threat but also against future 
variants that use the same technique.

Examples include:

\begin{itemize}
    \item encoded PowerShell execution
    \item unusual parent-child process chains
    \item suspicious authentication patterns
    \item rare outbound network connections
\end{itemize}

CTI provides the adversary context; detection engineering converts it into automated logic.

% --------------------------------------------------------------------
% SECTION 5: CTI-Driven Dashboard Design
% --------------------------------------------------------------------

\subsection*{CTI-Driven Dashboard Design}

CTI workflows also support dashboard creation. Intelligence about adversary behavior informs which panels should be 
included in dashboards and how they should be structured.

For example:

\begin{itemize}
    \item If a threat actor uses scheduled tasks for persistence → include a panel for scheduled task creation.
    \item If an actor relies on PowerShell → include a panel for PowerShell execution trends.
    \item If an actor uses rare outbound connections → include a panel for low-frequency network activity.
\end{itemize}

CTI ensures dashboards remain aligned with real-world threats rather than generic system activity.

% --------------------------------------------------------------------
% SECTION 6: Enrichment Workflows
% --------------------------------------------------------------------

\subsection*{Enrichment Workflows}

CTI also enhances investigations through enrichment. When analysts encounter suspicious events in Splunk or Sentinel, 
CTI provides context that helps determine whether the activity is benign or malicious.

Examples include:

\begin{itemize}
    \item enriching suspicious domains with threat actor associations
    \item enriching IP addresses with geolocation or reputation
    \item enriching file hashes with malware family information
\end{itemize}

Enrichment accelerates investigations and improves decision-making.

% --------------------------------------------------------------------
% SECTION 7: Feedback Loops
% --------------------------------------------------------------------

\subsection*{Feedback Loops}

CTI workflows rely on feedback loops. Threat hunters, SOC analysts, and detection engineers provide insight into:

\begin{itemize}
    \item which intelligence is useful
    \item which detections generate noise
    \item which dashboards need refinement
\end{itemize}

This feedback helps refine intelligence requirements and ensures that CTI remains aligned with operational needs. 
Over time, this iterative process strengthens the organization’s defensive posture and improves the quality of both 
intelligence and detections.

% --------------------------------------------------------------------
% SECTION 8: Why CTI Workflows Matter
% --------------------------------------------------------------------

\subsection*{Why CTI Workflows Matter}

The workflows described in this chapter form the backbone of operational CTI. By connecting intelligence to telemetry, 
hunts, detections, dashboards, and investigations, CTI becomes a practical force multiplier rather than a theoretical 
exercise. Later chapters expand these workflows into full playbooks and detection patterns tailored to specific adversary 
techniques.
